% !PRINTED 653
\weberpage{653}
%<<CONT
anwendet, und mit $A_1$, $A_2$, $A_3$, $A_4$ rationale Grössen bezeichnet:
\weq{14}{\begin{split}
K_0 &= A_1 + A_2 r + A_3 \varrho + A_4 r\varrho\\
K_1 &= A_1 - A_2 r + A_3 \varrho' - A_4 r\varrho'\\
K_2 &= A_1 + A_2 r - A_3 \varrho + A_4 r\varrho\\
% [sic] Print repeats $A_2$ as the coefficient of $\varrho'$ in the $K_3$
% row; the three parallel rows make $A_3$ the structurally expected coefficient.
K_3 &= A_1 - A_2 r - A_2 \varrho' - A_4 r\varrho'.
\end{split}}
Diese Grössen $K_0$, $K_1$, $K_2$, $K_3$ sind selbst wieder die Wurzeln einer cyklischen biquadratischen Gleichung. Sie haben nur scheinbar eine allgemeinere Form, als die $k$; denn setzt man $K_0$ in die Form
\[
K_0 = A_1 + A_2 r + \sqrt{\varrho^2(A_3+A_4 r)^2},
\]
so erkennt man die Form von $k_0$ in (9) wieder, natürlich mit veränderten $a$, $b$, $c$.
% [printed page ends here: a short horizontal rule follows below this sentence, marking the close of \S.\,188; the remainder of the printed page is blank]
